\clearpage
\fancyhead[L]{\small KOMET.}
\hypertarget{native-456-comet}{}
\begin{center}
{\large KOMET.}
\end{center}

\begin{center}
\rule{6cm}{0.4pt}\\
Göttingische gelehrte Anzeigen. Stück 57. Seite 561..565. 1826. April 10.\\
\rule{6cm}{0.4pt}
\end{center}

Unsere Zeitschrift hatte 1819 in den Nummern 28 und 83 [S. 417 und 420 dieses Bandes] als Erste öffentlich über die wichtige Entdeckung berichtet, durch die wir einen Kometen mit einer Umlaufzeit von nur drei Jahren kennengelernt haben. Er wurde nach dem Astronomen, der sich um seine Theorie so verdient gemacht hat, der \SourceLeaf{457}\textsc{Encke}sche Komet genannt. Der Bericht schloss mit dem Wunsch, dass \emph{diese höchst bemerkenswerte Entdeckung die Freunde der Astronomie anspornen möge, mit doppeltem Eifer nach Kometen zu suchen, da alles darauf hindeute, dass diese Entdeckung nur der Anfang einer unermesslichen, nach und nach reifenden Ernte sein werde.} Wie weit sich dieser \emph{Wunsch} erfüllt hat, zeigen die 14 neuen Kometen, die seitdem beobachtet und berechnet worden sind. Mit hoher Wahrscheinlichkeit können wir nun auch die erste Bestätigung der \emph{Vermutung} melden. Hauptmann \textsc{Biela} entdeckte am 28. Februar dieses Jahres in Josephsstadt in Böhmen im Sternbild Widder einen kleinen Kometen, der mit bloßem Auge unsichtbar war. Nach seiner Mitteilung wurde er auch an mehreren anderen Orten beobachtet. An unserer Sternwarte hat ihn Professor \textsc{Harding} bisher zweimal beobachtet. Einige Tage später entdeckte Herr \textsc{Gambard} in Marseille unabhängig davon denselben Kometen. Bei der großen Zahl neuer Kometen, die in den letzten Jahren gefunden wurden, würden wir diesen in unserer Zeitschrift nicht eigens erwähnen, wenn nicht schon die ersten Bahnberechnungen, die von mehreren verschiedenen Personen angestellt wurden und sich auch auf verschiedene Beobachtungen stützten, ein höchst bemerkenswertes Ergebnis erbracht hätten. Hofrat \textsc{Gauss} hat von fünf Astronomen die parabolischen Elemente erhalten, von Hauptmann von \textsc{Biela} selbst, Herrn \textsc{Clausen} in Altona, Professor \textsc{Encke} in Berlin, Dr.\ \textsc{Olbers} in Bremen und Professor \textsc{Schwerdt} in Speyer. Wir stellen sie hier zusammen:

\begin{center}
\begingroup\small\setlength{\tabcolsep}{5pt}
\begin{tabular}{rlll}
\multicolumn{4}{l}{Durchgang durch das Perihel} \\
1826. März & 15.45653 & v. B. & \\
 & 15.75422 & Cl. & \\
 & 18.2432 & Göttinger Z. & E. \\
 & 18.33773 & --- & O. \\
 & 17.620 & Sch. & \\
\end{tabular}
\par\medskip
\begin{tabular}{ll}
\multicolumn{2}{l}{Länge des Perihels} \\
$96^{\circ}\quad27'\quad33$ & v. B. \\
$95\quad48\quad15$ & Cl. \\
$107\quad55\quad31$ & E. \\
$112\quad39\quad24$ & O \\
$100\quad51\quad57$ & Sch. \\
\end{tabular}
\par\medskip
\begin{tabular}{ll}
\multicolumn{2}{l}{Aufsteigender Knoten} \\
$245^{\circ}\quad57'\quad0''$ & v. B. \\
$245\quad16\quad40$ & Cl. \\
$248\quad41\quad59$ & E. \\
$250\quad44\quad21$ & O. \\
$246\quad34\quad52$ & Sch. \\
\end{tabular}
\par\medskip
\begin{tabular}{ll}
\multicolumn{2}{l}{Neigung der Bahn} \\
$15^{\circ}\quad28'\quad23''$ & v. B. \\
$15\quad45\quad19$ & Cl. \\
$13\quad12\quad37$ & E. \\
$12\quad18\quad31$ & O. \\
$15\quad7\quad5$ & Sch. \\
\end{tabular}
\par\medskip
\begin{tabular}{ll}
\multicolumn{2}{l}{Logarithmus des kleinsten Abstands} \\
0.00506 & v. B. \\
0.0114369 & Cl. \\
9.96474 & E. \\
9.94460 & O. \\
9.99565 & Sch. \\
\end{tabular}
\endgroup
\end{center}
\noindent Rechtläufige Bewegung.

Angesichts der dürftigen Daten, auf denen diese Ergebnisse beruhen, sind die Unterschiede gering und werden sich bald ausgleichen lassen.

Sofort fällt auf, wie stark diese Elemente durchweg denen des zweiten Kometen von 1805 ähneln. Schon diese Ähnlichkeit allein macht es außerordentlich wahrscheinlich, dass beide Kometen identisch sind. Die Umstände beim Kometen von 1805 verstärken diese Wahrscheinlichkeit noch erheblich. Hofrat \textsc{Gauss}, der dessen Erscheinung sehr eingehend untersucht hatte, fand, dass die \SourceLeaf{458}Beobachtungen sehr auffällig auf eine kurze Umlaufzeit hinwiesen, obwohl sie nicht ausreichten, um diese aus der Erscheinung allein einigermaßen genau zu bestimmen. Am besten ließen sich die Beobachtungen mit einer Umlaufzeit von $4\frac{3}{4}$ Jahren vereinbaren. Nach jenen Untersuchungen hätte sich jedoch auch eine um ein Jahr kürzere oder um einige Jahre längere Umlaufzeit noch sehr gut und viel besser als eine parabolische Bewegung mit den Beobachtungen in Einklang bringen lassen. Zudem hatten die Elemente des Kometen von 1805 denen des Kometen von 1772 stark geähnelt. Das hatte zu der Vermutung geführt, dass beide identisch sein könnten. Der Unterschied war allerdings noch so beträchtlich, dass einige Astronomen die Identität für unwahrscheinlich hielten oder sogar sicher waren, dass es verschiedene Kometen waren. Hofrat \textsc{Gauss} zeigte aber schon damals, dass dieses Urteil nicht hinreichend begründet und eine vollständige Erklärung des Unterschieds keineswegs ausgeschlossen war. Obwohl er inzwischen nicht dazu gekommen ist, die umfangreichen Berechnungen anzustellen, die zur Entscheidung nötig wären, hielt er die Identität der beiden Kometen weiterhin für wahrscheinlich. Unter diesen Umständen ist es höchst bemerkenswert, dass die Zeit zwischen den Periheldurchgängen der Kometen von 1772 und 1805 und die Zeit zwischen den Durchgängen der Kometen von 1805 und 1826 annähernd im Verhältnis der Zahlen 5 und 3 stehen.

Wir dürfen daher mit hoher Wahrscheinlichkeit vermuten, dass alle drei Kometen identisch sind, dass dieser Komet von 1772 bis 1826 insgesamt acht Umläufe von durchschnittlich jeweils 6 Jahren und 9 Monaten vollzogen hat und dass die Unterschiede in den Elementen von 1772 und 1805, die einige Astronomen an der vermuteten Identität hatten zweifeln lassen, vor allem auf Störungen durch Jupiter zurückgingen, dem der Komet in der Zwischenzeit sehr nahe gekommen ist. Ohne diesen letzten Umstand wäre es vorerst noch ebenso wahrscheinlich, dass der Komet in der Zwischenzeit 10 und 6, also insgesamt 16 Umläufe von jeweils $3\frac{3}{8}$ Jahren vollzogen hätte.

Es ist sehr zu wünschen, dass der von Herrn von \textsc{Biela} entdeckte Komet, dessen Beobachtungen Mitte März durch den Mondschein unterbrochen wurden, nach dessen Ende irgendwo noch mehrmals gut beobachtet werden kann. Sollte dies nicht gelingen, so werden selbst die bereits gewonnenen Beobachtungen, sobald sie angemessen ausgewertet sind, ausreichen, um über die sich so unwiderstehlich aufdrängende Vermutung zu entscheiden, dass diese drei Kometen identisch sind.